% !TEX root = praca.tex

\chapter{Wprowadzenie}
\label{cha:wprowadzenie}

Coś o magazynach, sklepach, kompletacji zamówień. Coś o tym, że mamy magazyn, chcemy po nim jeździć aby zebrać zamówienie. (Np. Współczesne magazyny ...)

Definicja problemu komiwojażera (co to jest), dlaczego jest np trudny (skomplikowany obliczeniowo), dlaczego jest ważny w kontekście problemu kompletacji zamówień.

\section{Cel i zakres pracy}
\label{sec:celePracy}

Celem projektu dyplomowego jest zaprojektowanie i wdrożenie systemu wspierającego zarządzanie procesem automatycznej kompletacji zamówień w środowisku magazynowym. Projekt będzie oparty na integracji platformy mikrokomputerowej (np. Raspberry Pi, Arduino) jako jednostki sterującej pojazdem autonomicznym z aplikacją do zarządzania magazynem. Zaimplementowany zostanie system, który umożliwi pojazdowi zdalnie sterowanemu nawigację do wyznaczonego punktu kompletacji zamówienia. System może bazować na czujnikach odległości lub etykietach NFC. Kluczowe będzie zapewnienie stabilnej, dwukierunkowej komunikacji bezprzewodowej pomiędzy pojazdem a serwerem. Komunikacja będzie obejmować przesyłanie informacji o dotarciu do celu oraz odbieranie nowych zadań.
